%% toc: P3 39 39.3
\section{Proofs: finite-variable languages and pebble games}
\label{sec:annex-p03-w03-p03-056-pebble}

Notation is that of \cref{sec:p03-w03-p03-056-pebble}. Formulas of
\(\mathrm{L}^k_{\infty\omega}\) are well-founded formation trees, so
statements about all formulas are proved by induction on formation. A
position \((\alpha,\beta)\) has domain \(D\subseteq[k]\); its pairs are the
pairs \((\alpha(i),\beta(i))\), \(i\in D\). For \(D'\subseteq D\) write
\((\alpha,\beta)\restriction D'\) for
\((\alpha\restriction D',\beta\restriction D')\).

\begin{annexitem}{annex:p03-w03-p03-056-coincidence}{thm:p03-w03-p03-056-coincidence}{Only free labels matter}
\begin{proof}
Induction on \(\varphi\), for all pairs \(\alpha,\alpha'\) at once.
If \(\varphi\) is atomic, every variable \(v_i\) occurring in it has
\(i\in\operatorname{fr}(\varphi)\), so each term has the same value under
\(\alpha\) and \(\alpha'\), and the atomic clause gives the same truth
value. If \(\varphi=\neg\psi\), then
\(\operatorname{fr}(\psi)=\operatorname{fr}(\varphi)\) and the induction
hypothesis for \(\psi\) applies. If \(\varphi=\bigwedge\Phi\) or
\(\bigvee\Phi\), each \(\psi\in\Phi\) has
\(\operatorname{fr}(\psi)\subseteq\operatorname{fr}(\varphi)\), so
\(\alpha,\alpha'\) are admissible for \(\psi\) and agree on its free labels;
by the induction hypothesis every member of \(\Phi\) has the same truth value
under both, hence so does \(\varphi\). If \(\varphi=\exists v_i\psi\), let
\(a\in A\). The assignments \(\alpha[i\mapsto a]\) and
\(\alpha'[i\mapsto a]\) are admissible for \(\psi\), since
\(\operatorname{fr}(\psi)\subseteq\operatorname{fr}(\varphi)\cup\{i\}\), and
they agree on \(\operatorname{fr}(\psi)\): at \(i\) both give \(a\), and at
\(j\in\operatorname{fr}(\psi)\setminus\{i\}\subseteq\operatorname{fr}(\varphi)\)
they give \(\alpha(j)=\alpha'(j)\). By the induction hypothesis
\(\mathbf A\models\psi[\alpha[i\mapsto a]]\iff
\mathbf A\models\psi[\alpha'[i\mapsto a]]\) for every \(a\), so the
existential clause gives the same value; the universal case is the same
with ``for every \(a\)'' in place of ``for some \(a\)''.
\end{proof}
In particular, if \((\mathbf A,\alpha)\) and \((\mathbf B,\beta)\) agree on
all formulas of some class with free labels in \(D\), and \(D'\subseteq D\),
then \((\mathbf A,\alpha\restriction D')\) and
\((\mathbf B,\beta\restriction D')\) agree on the formulas of that class with
free labels in \(D'\): such a formula has the same value under
\(\alpha\restriction D'\) as under \(\alpha\), and likewise in \(\mathbf B\).
This \emph{restriction property} is used below for removal moves.
\end{annexitem}

\begin{annexitem}{annex:p03-w03-p03-056-renaming}{thm:p03-w03-p03-056-renaming}{Counting variables versus fixing them}
\begin{proof}
Induction on \(\varphi\). Every variable of \(\varphi^\pi\) lies in
\(\pi[W]\subseteq\{v_1,\ldots,v_k\}\), and the clauses of
\cref{def:p03-w03-p03-039-language} for \(\neg\), \(\bigwedge\),
\(\bigvee\), \(\exists\), \(\forall\) are matched by (G2)--(G4); for a set
\(\Phi\), \((\bigwedge\Phi)^\pi=\bigwedge\{\psi^\pi:\psi\in\Phi\}\) is the
conjunction of a set (the image of \(\Phi\)), and likewise for
\(\bigvee\). So \(\varphi^\pi\in\mathrm{L}^k_{\infty\omega}(\sigma)\). Since
\(\pi\) is injective, \(s\circ\pi^{-1}\) is a well-defined assignment on
\(\pi[S]\), and it is admissible for \(\varphi^\pi\) because the free
variables of \(\varphi^\pi\) are the images of those of \(\varphi\).

Atomic \(\varphi\): a variable \(w\) occurring in \(\varphi\) is free, so
\(w\in S\), and its value \(s(w)\) equals the value
\((s\circ\pi^{-1})(\pi(w))\) of \(\pi(w)\); constants are unchanged. The
atomic clause gives the same truth value. Negation, \(\bigwedge\) and
\(\bigvee\): each immediate subformula has free variables among those of
\(\varphi\), hence in \(S\), and the induction hypothesis applies to each.
Quantifier: \(\varphi=\exists w\,\psi\), so
\(\varphi^\pi=\exists\pi(w)\,\psi^\pi\). For \(a\in A\), the assignment
\(s[w\mapsto a]\) has domain \(S\cup\{w\}\subseteq W\), which contains the
free variables of \(\psi\). Moreover
\((s[w\mapsto a])\circ\pi^{-1}=(s\circ\pi^{-1})[\pi(w)\mapsto a]\): both
send \(\pi(w)\) to \(a\), and for \(u\in S\setminus\{w\}\) both send
\(\pi(u)\) to \(s(u)\), because \(\pi(u)\neq\pi(w)\) by injectivity. Hence,
by the induction hypothesis,
\(\mathbf A\models\psi[s[w\mapsto a]]\iff
\mathbf A\models\psi^\pi[(s\circ\pi^{-1})[\pi(w)\mapsto a]]\) for every
\(a\), and the existential clause gives
\(\mathbf A\models\varphi[s]\iff\mathbf A\models\varphi^\pi[s\circ\pi^{-1}]\).
The universal case is identical with ``every'' for ``some''.
\end{proof}
\end{annexitem}

\begin{annexitem}{annex:p03-w03-p03-056-walks}{thm:p03-w03-p03-056-walks}{Walks of every length with three variables}
\begin{proof}
\emph{Syntax.} \(\theta_1\) is atomic with free labels \(\{1,2\}\) and rank
\(0\). If \(\operatorname{fr}(\theta_n)=\{1,2\}\) and
\(\operatorname{qr}(\theta_n)=2(n-1)\), then
\(\exists v_1(v_3=v_1\land\theta_n)\) has free labels
\((\{1,3\}\cup\{1,2\})\setminus\{1\}=\{2,3\}\) and rank \(2(n-1)+1\);
conjoining \(E(v_1,v_3)\) gives free labels \(\{1,2,3\}\) and the same rank;
and \(\exists v_3\) gives free labels \(\{1,2\}\) and rank \(2n\). All
conjunctions are binary and only \(v_1,v_2,v_3\) occur, so each
\(\theta_n\) is an \(\mathrm{FO}^3_{\omega\omega}\)-formula.

\emph{Semantics.} Say that \(e_0,\ldots,e_n\) is a walk of length \(n\) from
\(e_0\) to \(e_n\) when \((e_{j-1},e_j)\in E^{\mathbf A}\) for
\(1\leq j\leq n\). We prove, by induction on \(n\geq1\), that for every
assignment \(\alpha\) admissible for \(\theta_n\),
\(\mathbf A\models\theta_n[\alpha]\) if and only if there is a walk of length
\(n\) from \(\alpha(1)\) to \(\alpha(2)\). For \(n=1\) this is the atomic
clause. For the step,
\(\mathbf A\models\theta_{n+1}[\alpha]\) holds if and only if there is
\(c\in A\) with \((\alpha(1),c)\in E^{\mathbf A}\) and
\(\mathbf A\models\exists v_1(v_3=v_1\land\theta_n)[\alpha[3\mapsto c]]\);
the latter holds if and only if there is \(d\in A\) with \(c=d\) and
\(\mathbf A\models\theta_n[\gamma]\), where
\(\gamma=\alpha[3\mapsto c][1\mapsto d]\). So \(d=c\), and \(\gamma\)
sends \(1\) to \(c\) and \(2\) to \(\alpha(2)\). By the induction hypothesis
\(\mathbf A\models\theta_n[\gamma]\) if and only if there is a walk of
length \(n\) from \(c\) to \(\alpha(2)\). Thus
\(\mathbf A\models\theta_{n+1}[\alpha]\) if and only if some \(c\) is an
\(E\)-successor of \(\alpha(1)\) from which a walk of length \(n\) leads to
\(\alpha(2)\). Prefixing \(\alpha(1)\) to such a walk gives a walk of length
\(n+1\) from \(\alpha(1)\) to \(\alpha(2)\); conversely the second vertex of
a walk of length \(n+1\) is such a \(c\).

\emph{The sentence.} It is built by (G1)--(G4) with variables
\(v_1,v_2,v_3\), so it lies in \(\mathrm{L}^3_{\infty\omega}\), and it has
no free labels. It holds in \(\mathbf A\) if and only if for all
\(a,b\in A\) either \(a=b\) or, for some \(n\geq1\), there is a walk of
length \(n\) from \(a\) to \(b\); that is, every vertex reaches every other
vertex. When \(E^{\mathbf A}\) is symmetric this says that any two distinct
vertices are joined by a walk, and a shortest such walk repeats no vertex,
so it is a path; this is connectivity.
\end{proof}
\end{annexitem}

\begin{annexitem}{annex:p03-w03-p03-056-finitely-many}{thm:p03-w03-p03-056-finitely-many}{Finitely many first-order types of bounded rank}
\begin{proof}
\emph{Step 1: Boolean normal forms.} Let \(G\) be a finite set of
\(\mathrm{FO}^k_{\omega\omega}\)-formulas with free labels in \(D\). For
\(e:G\to\{0,1\}\) put \(\chi_e=\bigwedge\{g^{e(g)}:g\in G\}\), where
\(g^1=g\) and \(g^0=\neg g\); for \(S\subseteq\{0,1\}^G\) put
\(\delta_S=\bigvee\{\chi_e:e\in S\}\). There are \(2^{2^{\lvert G\rvert}}\)
formulas \(\delta_S\); each is in \(\mathrm{FO}^k_{\omega\omega}\) (all index
sets are finite), has free labels in \(D\), and has rank at most
\(\max\{\operatorname{qr}(g):g\in G\}\) (rank \(0\) if \(G=\emptyset\)).
Call a formula a \emph{Boolean combination over \(G\)} if it is built from
members of \(G\) by \(\neg\) and finite \(\bigwedge,\bigvee\). For a
structure \(\mathbf A\) and \(\alpha:D\to A\), let
\(e_{\mathbf A,\alpha}(g)=1\) if \(\mathbf A\models g[\alpha]\) and \(0\)
otherwise. By induction on the construction of a Boolean combination
\(\psi\) over \(G\) there is a function \(f_\psi:\{0,1\}^G\to\{0,1\}\) with
\(\mathbf A\models\psi[\alpha]\iff f_\psi(e_{\mathbf A,\alpha})=1\) for all
\(\mathbf A,\alpha\): for \(\psi=g\in G\) take
\(f_\psi(e)=e(g)\); negation, finite conjunction and finite disjunction are
computed pointwise from the functions of the parts. Put
\(S_\psi=\{e:f_\psi(e)=1\}\). Since \(\mathbf A\models\chi_e[\alpha]\)
exactly when \(e=e_{\mathbf A,\alpha}\), we get
\(\mathbf A\models\delta_{S_\psi}[\alpha]\iff e_{\mathbf A,\alpha}\in S_\psi
\iff\mathbf A\models\psi[\alpha]\). So every Boolean combination over \(G\)
is equivalent over \(D\) to one of the finitely many \(\delta_S\).

\emph{Step 2: atoms.} Let \(\mathrm{At}_D\) be the set of atomic formulas
whose terms lie in \(T_D=\{v_i:i\in D\}\cup\{c_1,\ldots,c_s\}\). Since
\(\lvert T_D\rvert\leq k+s\) and \(\sigma\) has finitely many relation
symbols, \(\lvert\mathrm{At}_D\rvert\leq(k+s)^2+\sum_{R\in\sigma}
(k+s)^{\operatorname{arity}(R)}\) is finite. An atomic formula has free
labels in \(D\) exactly when it belongs to \(\mathrm{At}_D\).

\emph{Step 3: induction on \(m\).} For \(m=0\): an
\(\mathrm{FO}^k_{\omega\omega}\)-formula of rank \(0\) contains no
quantifier, so, by induction on its formation, it is a Boolean combination
of its atomic subformulas; each of these has free labels among those of the
formula, hence in \(D\), and so lies in \(\mathrm{At}_D\). Step 1 with
\(G=\mathrm{At}_D\) gives \(F_{0,D}=\{\delta_S:S\subseteq\{0,1\}^{\mathrm{At}_D}\}\).

Suppose \(F_{m,D'}\) has been found for every \(D'\subseteq[k]\). Let
\[
 G_{m+1,D}=\mathrm{At}_D\cup\{\exists v_i f,\ \forall v_i f:
   i\in[k],\ f\in F_{m,D\cup\{i\}}\},
\]
a finite set of \(\mathrm{FO}^k_{\omega\omega}\)-formulas of rank at most
\(m+1\) with free labels in \(D\) (since
\(\operatorname{fr}(f)\subseteq D\cup\{i\}\)). We show by induction on
formation that every \(\mathrm{FO}^k_{\omega\omega}\)-formula \(\psi\) with
\(\operatorname{qr}(\psi)\leq m+1\) and \(\operatorname{fr}(\psi)\subseteq D\)
is equivalent over \(D\) to a Boolean combination over \(G_{m+1,D}\).
Atomic \(\psi\) lie in \(\mathrm{At}_D\). If \(\psi\) is \(\neg\psi'\) or a
finite \(\bigwedge\) or \(\bigvee\), its immediate subformulas have rank at
most \(m+1\) and free labels in \(D\); replacing each by an equivalent
Boolean combination over \(G_{m+1,D}\) gives an equivalent Boolean
combination, since the truth clauses for \(\neg,\bigwedge,\bigvee\) depend
only on the truth values of the parts. If \(\psi=\exists v_i\psi'\), then
\(\operatorname{qr}(\psi')\leq m\) and
\(\operatorname{fr}(\psi')\subseteq D\cup\{i\}\), so \(\psi'\) is equivalent
over \(D\cup\{i\}\) to some \(f\in F_{m,D\cup\{i\}}\). For any
\(\mathbf A\) and \(\alpha:D\to A\) and any \(a\in A\), the assignment
\(\alpha[i\mapsto a]\) has domain \(D\cup\{i\}\), so
\(\mathbf A\models\psi'[\alpha[i\mapsto a]]\iff
\mathbf A\models f[\alpha[i\mapsto a]]\); hence \(\psi\) is equivalent over
\(D\) to \(\exists v_if\in G_{m+1,D}\). The universal case gives
\(\forall v_if\). Step 1 with \(G=G_{m+1,D}\) now gives
\(F_{m+1,D}=\{\delta_S:S\subseteq\{0,1\}^{G_{m+1,D}}\}\), whose members have
rank at most \(m+1\) and free labels in \(D\).
\end{proof}
\end{annexitem}

\begin{annexitem}{annex:p03-w03-p03-056-atomic}{thm:p03-w03-p03-056-atomic}{Good positions are atomic agreement}
\begin{proof}
Every pair in \(\rho=\rho_{\alpha,\beta}\) has the form
\((t^{\mathbf A}[\alpha],t^{\mathbf B}[\beta])\) for a term \(t\in T_D\)
(\(t=v_i\) for the pair of label \(i\), \(t=c\) for a constant), and
conversely every \(t\in T_D\) yields such a pair.

Suppose \(\rho\) is a partial isomorphism \(h\). For \(t\in T_D\),
\(t^{\mathbf A}[\alpha]\in\operatorname{dom}(h)\) and
\(h(t^{\mathbf A}[\alpha])=t^{\mathbf B}[\beta]\), because \(h\) is a function
containing that pair. For an atomic \(t=t'\) with terms in \(T_D\):
\(t^{\mathbf A}[\alpha]=t'^{\mathbf A}[\alpha]\) if and only if their
\(h\)-images are equal, by injectivity, that is, if and only if
\(t^{\mathbf B}[\beta]=t'^{\mathbf B}[\beta]\). For \(R(t_1,\ldots,t_r)\): the
values \(t_j^{\mathbf A}[\alpha]\) lie in \(\operatorname{dom}(h)\), and
\(h\) preserves and reflects \(R\), so
\(R^{\mathbf A}(t_1^{\mathbf A}[\alpha],\ldots)\iff
R^{\mathbf B}(t_1^{\mathbf B}[\beta],\ldots)\).

Conversely suppose agreement on all atomic formulas with free labels in
\(D\). \emph{Function:} if \((a,b),(a,b')\in\rho\), choose terms
\(t,t'\in T_D\) with \(a=t^{\mathbf A}[\alpha]=t'^{\mathbf A}[\alpha]\),
\(b=t^{\mathbf B}[\beta]\) and \(b'=t'^{\mathbf B}[\beta]\). The atomic
formula \(t=t'\) holds in \(\mathbf A\) under \(\alpha\), hence in
\(\mathbf B\) under \(\beta\), so \(b=b'\). \emph{Injective:} the same
argument with the roles of \(\mathbf A\) and \(\mathbf B\) exchanged.
\emph{Constants:} \((c^{\mathbf A},c^{\mathbf B})\in\rho\) by definition, so
the domain contains every \(c^{\mathbf A}\) and \(h(c^{\mathbf A})=c^{\mathbf B}\).
\emph{Relations:} for \(a_1,\ldots,a_r\in\operatorname{dom}(h)\) choose
\(t_j\in T_D\) with \(a_j=t_j^{\mathbf A}[\alpha]\); then
\(h(a_j)=t_j^{\mathbf B}[\beta]\), and agreement on
\(R(t_1,\ldots,t_r)\) gives
\(R^{\mathbf A}(a_1,\ldots,a_r)\iff R^{\mathbf B}(h(a_1),\ldots,h(a_r))\).
So \(h\) is a partial isomorphism in the sense of
\cref{def:p03-w03-p03-056-vocabulary}.
\end{proof}
\end{annexitem}

\begin{annexitem}{annex:p03-w03-p03-056-finite-round}{thm:p03-w03-p03-056-finite-round}{Finite-round pebble games and bounded-rank \(k\)-variable formulas}
\begin{proof}
\emph{(1)\(\Rightarrow\)(2).} Let \(\Sigma\) be winning in
\(\mathrm{PG}^n_k(\mathbf A,\mathbf B)\) from \((\alpha,\beta)\). A sequence
of Spoiler moves without removals is a sequence of Spoiler moves of the full
game, so \(\Sigma\) answers it, and every play of the restricted game that
follows \(\Sigma\) is a play of the full game that follows \(\Sigma\); all
its positions are good.

\emph{(2)\(\Rightarrow\)(3).} We prove by induction on the formation of
\(\varphi\in\mathrm{L}^k_{\infty\omega}\): for every \(m\in\N\) with
\(\operatorname{qr}(\varphi)\leq m\) and every position \((\alpha,\beta)\)
whose domain contains \(\operatorname{fr}(\varphi)\) and from which the
Duplicator has a winning strategy in the \(m\)-round game without removals,
\(\mathbf A\models\varphi[\alpha]\iff\mathbf B\models\varphi[\beta]\). Item~3
is the case \(m=n\).

\emph{Atomic \(\varphi\).} A winning strategy makes the initial position
good, so by \cref{thm:p03-w03-p03-056-atomic} the position agrees on
\(\varphi\).
\emph{\(\neg\), \(\bigwedge\), \(\bigvee\).} The immediate subformulas have
rank at most \(\operatorname{qr}(\varphi)\leq m\) and free labels in the
domain, so the induction hypothesis applies to each, at the same \(m\) and
position, and the truth clauses transfer agreement.
\emph{\(\varphi=\exists v_i\psi\).} Then
\(\operatorname{qr}(\psi)+1\leq m\), so \(m\geq1\). Suppose
\(\mathbf A\models\varphi[\alpha]\) and choose \(a\) with
\(\mathbf A\models\psi[\alpha[i\mapsto a]]\). Let the Spoiler's first move
place the \(\mathbf A\)-pebble \(i\) on \(a\), and let \(b\) be the answer of
the winning strategy \(\Sigma\). The new position
\((\alpha[i\mapsto a],\beta[i\mapsto b])\) has domain containing
\(\operatorname{fr}(\psi)\subseteq\operatorname{fr}(\varphi)\cup\{i\}\), and
by the residual-strategy remark in \cref{def:p03-w03-p03-056-pebble-game}
the Duplicator has a winning strategy from it in the \((m-1)\)-round game
without removals. Since \(\operatorname{qr}(\psi)\leq m-1\), the induction
hypothesis gives \(\mathbf B\models\psi[\beta[i\mapsto b]]\), hence
\(\mathbf B\models\varphi[\beta]\). If \(\mathbf B\models\varphi[\beta]\),
choose \(b\) with \(\mathbf B\models\psi[\beta[i\mapsto b]]\); the Spoiler
places the \(\mathbf B\)-pebble \(i\) on \(b\), the strategy answers some
\(a\in A\), the position \((\alpha[i\mapsto a],\beta[i\mapsto b])\) is
winning for \(m-1\) rounds, and the induction hypothesis gives
\(\mathbf A\models\psi[\alpha[i\mapsto a]]\), so
\(\mathbf A\models\varphi[\alpha]\).
\emph{\(\varphi=\forall v_i\psi\).} If \(\mathbf A\not\models\varphi[\alpha]\),
choose \(a\) with \(\mathbf A\not\models\psi[\alpha[i\mapsto a]]\); the
Spoiler places the \(\mathbf A\)-pebble \(i\) on \(a\), the answer \(b\)
leads to a position winning for \(m-1\) rounds, and the induction hypothesis
gives \(\mathbf B\not\models\psi[\beta[i\mapsto b]]\), so
\(\mathbf B\not\models\varphi[\beta]\). If
\(\mathbf B\not\models\varphi[\beta]\), choose \(b\) with
\(\mathbf B\not\models\psi[\beta[i\mapsto b]]\); the Spoiler places the
\(\mathbf B\)-pebble \(i\) on \(b\), the answer \(a\) leads to a position
winning for \(m-1\) rounds, and the induction hypothesis gives
\(\mathbf A\not\models\psi[\alpha[i\mapsto a]]\), so
\(\mathbf A\not\models\varphi[\alpha]\).

\emph{(3)\(\Rightarrow\)(4).} Every \(\mathrm{FO}^k_{\omega\omega}\)-formula is
an \(\mathrm{L}^k_{\infty\omega}\)-formula, with the same rank.

\emph{(4)\(\Rightarrow\)(1).} For \(m\in\N\) let \(J_m\) be the set of
positions \((\alpha',\beta')\), with any domain \(D'\subseteq[k]\), such that
\((\mathbf A,\alpha')\equiv^{k,m}_{\omega\omega}(\mathbf B,\beta')\). We
establish four properties.
\begin{enumerate}
\item[(a)] Every member of \(J_m\) is good: atomic formulas are
      \(\mathrm{FO}^k_{\omega\omega}\)-formulas of rank \(0\leq m\); apply
      \cref{thm:p03-w03-p03-056-atomic}.
\item[(b)] \(J_m\subseteq J_{m-1}\) for \(m\geq1\): fewer formulas are
      required to agree.
\item[(c)] If \((\alpha',\beta')\in J_m\) has domain \(D'\) and
      \(D''\subseteq D'\), then \((\alpha',\beta')\restriction D''\in J_m\),
      by the restriction property after
      \cref{annex:p03-w03-p03-056-coincidence}.
\item[(d)] If \((\alpha',\beta')\in J_m\), \(m\geq1\), \(i\in[k]\) and
      \(a\in A\), then there is \(b\in B\) with
      \((\alpha'[i\mapsto a],\beta'[i\mapsto b])\in J_{m-1}\); and
      symmetrically for \(b\in B\).
\end{enumerate}
Proof of (d). Let \(D'\) be the domain and \(E=D'\cup\{i\}\). Let \(\Psi\) be
the set of members of \(F_{m-1,E}\) (from
\cref{thm:p03-w03-p03-056-finitely-many}) that hold in \(\mathbf A\) under
\(\alpha'[i\mapsto a]\), together with the formula \(v_i=v_i\); it is a
finite set. Put \(\chi=\bigwedge\Psi\). Then \(\exists v_i\chi\) is an
\(\mathrm{FO}^k_{\omega\omega}\)-formula of rank at most \(m\) with free labels
in \(D'\), and \(\mathbf A\models\exists v_i\chi[\alpha']\), witnessed by
\(a\). As \((\alpha',\beta')\in J_m\),
\(\mathbf B\models\exists v_i\chi[\beta']\); choose \(b\) with
\(\mathbf B\models\chi[\beta'[i\mapsto b]]\). Now let \(\varphi\) be any
\(\mathrm{FO}^k_{\omega\omega}\)-formula of rank at most \(m-1\) with free
labels in \(E\). If \(\mathbf A\models\varphi[\alpha'[i\mapsto a]]\), choose
\(f\in F_{m-1,E}\) equivalent to \(\varphi\) over \(E\); then \(f\) holds in
\(\mathbf A\) under \(\alpha'[i\mapsto a]\), so \(f\in\Psi\), so
\(\mathbf B\models f[\beta'[i\mapsto b]]\), and by equivalence
\(\mathbf B\models\varphi[\beta'[i\mapsto b]]\). If
\(\mathbf A\not\models\varphi[\alpha'[i\mapsto a]]\), apply the same
argument to \(\neg\varphi\), which has the same rank and free labels; it
gives \(\mathbf B\not\models\varphi[\beta'[i\mapsto b]]\). Hence the new
position is in \(J_{m-1}\). For \(b\in B\), let \(\Psi'\) consist of
\(v_i=v_i\) and the members of \(F_{m-1,E}\) true in \(\mathbf B\) under
\(\beta'[i\mapsto b]\), and \(\chi'=\bigwedge\Psi'\). Then
\(\mathbf B\models\exists v_i\chi'[\beta']\), so
\(\mathbf A\models\exists v_i\chi'[\alpha']\); a witness \(a\) satisfies
every member of \(\Psi'\) under \(\alpha'[i\mapsto a]\), and the argument
just given, with \(\mathbf A\) and \(\mathbf B\) exchanged, shows that
\((\alpha'[i\mapsto a],\beta'[i\mapsto b])\in J_{m-1}\).

\emph{The strategy.} By hypothesis \((\alpha,\beta)\in J_n\). The Duplicator
keeps the invariant that after \(j\) rounds (\(0\leq j\leq n\)) the position
lies in \(J_{n-j}\). It holds for \(j=0\). If it holds after \(j<n\) rounds,
then in round \(j+1\): on a placement of label \(i\) on \(a\in A\) (or
\(b\in B\)) she answers with the element given by (d) for
\(m=n-j\geq1\), and the new position lies in \(J_{n-j-1}\); on a removal the
new position is a restriction, which lies in \(J_{n-j}\) by (c), hence in
\(J_{n-j-1}\) by (b). Every position of the play therefore lies in some
\(J_m\) with \(m\geq0\) and is good by (a). This strategy is winning.
\end{proof}
\end{annexitem}

\begin{annexitem}{annex:p03-w03-p03-056-eternal}{thm:p03-w03-p03-056-eternal}{Eternal pebble games and \(\mathrm{L}^k_{\infty\omega}\)-equivalence}
\begin{proof}
\emph{(1)\(\Rightarrow\)(2) and (1)\(\Rightarrow\)(3).} In each variant the
Spoiler's available moves are moves of \(\mathrm{PG}^\infty_k\): placement of
any pebble, respectively placement of a pebble not on the board, and
removal. A winning strategy for the full game, applied to the variant's
Spoiler sequences, produces only plays of the full game that follow it, so
all their positions are good.

\emph{(2)\(\Rightarrow\)(4) and (3)\(\Rightarrow\)(4).} Fix one of the two
variants and call a position \emph{winning} if the Duplicator has a winning
strategy from it in that eternal variant. By the residual-strategy remark,
after any Spoiler move from a winning position and the strategy's answer,
the new position is again winning. We prove by induction on
\(\varphi\in\mathrm{L}^k_{\infty\omega}\): at every winning position whose
domain contains \(\operatorname{fr}(\varphi)\),
\(\mathbf A\models\varphi[\alpha]\iff\mathbf B\models\varphi[\beta]\).
Atomic formulas and the Boolean clauses are handled exactly as in
\cref{annex:p03-w03-p03-056-finite-round}, without the rank bookkeeping.
Let \(\varphi=\exists v_i\psi\), let \((\alpha,\beta)\) be winning with
domain \(D\supseteq\operatorname{fr}(\varphi)\), and suppose
\(\mathbf A\models\psi[\alpha[i\mapsto a]]\).
\begin{itemize}
\item If the variant is the one without removals, or if \(i\notin D\), the
      Spoiler places the \(\mathbf A\)-pebble \(i\) on \(a\) (permitted:
      in the second variant pebble \(i\) is not on the board). The answer
      \(b\) yields the winning position
      \((\alpha[i\mapsto a],\beta[i\mapsto b])\).
\item If the variant is Kolaitis's and \(i\in D\), the Spoiler first removes
      pebble \(i\), giving the winning position
      \((\alpha,\beta)\restriction(D\setminus\{i\})\), and then places the
      now unused \(\mathbf A\)-pebble \(i\) on \(a\); the answer \(b\)
      yields the winning position
      \(((\alpha\restriction(D\setminus\{i\}))[i\mapsto a],
      (\beta\restriction(D\setminus\{i\}))[i\mapsto b])\), which is
      \((\alpha[i\mapsto a],\beta[i\mapsto b])\).
\end{itemize}
In both cases the induction hypothesis for \(\psi\) at this winning position
(whose domain contains \(\operatorname{fr}(\psi)\subseteq D\cup\{i\}\)) gives
\(\mathbf B\models\psi[\beta[i\mapsto b]]\), so
\(\mathbf B\models\varphi[\beta]\).

The same two bullet moves, made with the \(\mathbf B\)-pebble \(i\) on an
element \(b\in B\), produce an answer \(a\in A\) and the winning position
\((\alpha[i\mapsto a],\beta[i\mapsto b])\). Applied to a witness \(b\) of
\(\mathbf B\models\varphi[\beta]\), the induction hypothesis gives
\(\mathbf A\models\psi[\alpha[i\mapsto a]]\), hence
\(\mathbf A\models\varphi[\alpha]\). For \(\varphi=\forall v_i\psi\): if
\(\mathbf A\not\models\psi[\alpha[i\mapsto a]]\) for some \(a\), the moves
on \(a\) give \(b\) with \(\mathbf B\not\models\psi[\beta[i\mapsto b]]\) by
the induction hypothesis, so \(\mathbf B\not\models\varphi[\beta]\); and if
\(\mathbf B\not\models\psi[\beta[i\mapsto b]]\) for some \(b\), the moves on
\(b\) give \(a\) with \(\mathbf A\not\models\psi[\alpha[i\mapsto a]]\), so
\(\mathbf A\not\models\varphi[\alpha]\).

\emph{(4)\(\Rightarrow\)(1).} Let \(I\) be the set of positions
\((\alpha',\beta')\), with any domain \(D'\subseteq[k]\), such that
\((\mathbf A,\alpha')\equiv^k_{\infty\omega}(\mathbf B,\beta')\).
\begin{enumerate}
\item[(a)] Members of \(I\) are good, by \cref{thm:p03-w03-p03-056-atomic}.
\item[(b)] \(I\) is closed under restriction, by the restriction property.
\item[(c)] If \((\alpha',\beta')\in I\) has domain \(D'\), \(i\in[k]\) and
      \(a\in A\), then some \(b\in B\) has
      \((\alpha'[i\mapsto a],\beta'[i\mapsto b])\in I\); symmetrically for
      \(b\in B\).
\end{enumerate}
Proof of (c). Suppose not. Then for each \(b\in B\) there is an
\(\mathrm{L}^k_{\infty\omega}\)-formula with free labels in \(D'\cup\{i\}\)
on which \((\mathbf A,\alpha'[i\mapsto a])\) and
\((\mathbf B,\beta'[i\mapsto b])\) disagree; replacing it by its negation if
necessary, and using the axiom of choice, choose for each \(b\) such a
formula \(\psi_b\) with \(\mathbf A\models\psi_b[\alpha'[i\mapsto a]]\) and
\(\mathbf B\not\models\psi_b[\beta'[i\mapsto b]]\). The collection
\(\Phi=\{\psi_b:b\in B\}\) is a set, the image of \(B\), so
\(\vartheta=\exists v_i\bigwedge\Phi\) is an
\(\mathrm{L}^k_{\infty\omega}\)-formula, and its free labels lie in \(D'\).
It holds in \(\mathbf A\) under \(\alpha'\), witnessed by \(a\). If it held in
\(\mathbf B\) under \(\beta'\), some \(b'\) would satisfy every \(\psi_b\)
under \(\beta'[i\mapsto b']\), in particular \(\psi_{b'}\), contrary to the
choice of \(\psi_{b'}\). So \(\vartheta\) separates \((\mathbf A,\alpha')\)
from \((\mathbf B,\beta')\), contradicting \((\alpha',\beta')\in I\). For
\(b\in B\) the argument is the same with the structures exchanged: if no
\(a\in A\) works, choose for each \(a\) a formula \(\psi'_a\) true in
\(\mathbf B\) under \(\beta'[i\mapsto b]\) and false in \(\mathbf A\) under
\(\alpha'[i\mapsto a]\); then \(\exists v_i\bigwedge\{\psi'_a:a\in A\}\)
holds in \(\mathbf B\) under \(\beta'\) and fails in \(\mathbf A\) under
\(\alpha'\).

The Duplicator's strategy is to answer every placement by an element given
by (c) and to answer removals as the rules require. By induction on the
number of rounds, every position of the play is in \(I\): the initial one by
hypothesis, placements by (c), removals by (b). By (a) all positions are good,
so the strategy wins every play; it also answers every placement of a
pebble not on the board and needs no removal, so it wins the two variants as
well.

\emph{(4)\(\Rightarrow\)\(\equiv^k_{\omega\omega}\).}
\(\mathrm{FO}^k_{\omega\omega}\subseteq\mathrm{L}^k_{\infty\omega}\).

\emph{Finite structures.} Let \(\mathbf A,\mathbf B\) be finite and
\((\mathbf A,\alpha)\equiv^k_{\omega\omega}(\mathbf B,\beta)\). Let
\(I^{\mathrm{FO}}\) be the set of positions with
\(\equiv^k_{\omega\omega}\). Properties (a) and (b) hold as before, since
atomic formulas are first-order. For (c), repeat the argument with each
\(\psi_b\) chosen in \(\mathrm{FO}^k_{\omega\omega}\); because \(B\) is
finite, \(\Phi\) is finite and \(\vartheta=\exists v_i\bigwedge\Phi\) is an
\(\mathrm{FO}^k_{\omega\omega}\)-formula, so the contradiction is reached
inside \(\mathrm{FO}^k_{\omega\omega}\). For a placement in \(\mathbf B\)
the conjunction \(\bigwedge\{\psi'_a:a\in A\}\) is finite because \(A\) is
finite. The strategy that stays in \(I^{\mathrm{FO}}\) wins
\(\mathrm{PG}^\infty_k\), which is item~1.
\end{proof}
\end{annexitem}

\begin{annexitem}{annex:p03-w03-p03-056-monotone}{thm:p03-w03-p03-056-monotone}{Pebble numbers and the equivalence relation}
\begin{proof}
(1) Suppose the Duplicator wins \(\mathrm{PG}^\infty_{k'}(\mathbf A,\mathbf B)\)
from the empty position. By \cref{thm:p03-w03-p03-056-eternal} (items~1
and~4, with \(k'\) and \(D=\emptyset\)),
\(\mathbf A\equiv^{k'}_{\infty\omega}\mathbf B\). Every
\(\mathrm{L}^k_{\infty\omega}\)-sentence is an
\(\mathrm{L}^{k'}_{\infty\omega}\)-sentence
(\cref{def:p03-w03-p03-056-grammar}), so
\(\mathbf A\equiv^k_{\infty\omega}\mathbf B\), and the same theorem with
\(k\) gives a winning strategy in \(\mathrm{PG}^\infty_k(\mathbf A,\mathbf B)\).
For \(n\) rounds argue in the same way with
\cref{thm:p03-w03-p03-056-finite-round} (items~1 and~3): an
\(\mathrm{L}^k_{\infty\omega}\)-sentence of rank at most \(n\) is an
\(\mathrm{L}^{k'}_{\infty\omega}\)-sentence of the same rank, since the rank
is defined by the same clauses in both languages. The statement about the
Spoiler is the contrapositive.

(2) By \cref{thm:p03-w03-p03-056-eternal}, the Duplicator wins
\(\mathrm{PG}^\infty_k(\mathbf A,\mathbf B)\) exactly when, for every
\(\mathrm{L}^k_{\infty\omega}\)-sentence \(\varphi\),
\(\mathbf A\models\varphi\iff\mathbf B\models\varphi\). A relation of the
form ``\(f(\mathbf A)=f(\mathbf B)\)'', where \(f(\mathbf A)\) is the class
of \(\mathrm{L}^k_{\infty\omega}\)-sentences true in \(\mathbf A\), is
reflexive, symmetric and transitive, because equality is. For \(n\) rounds
use \cref{thm:p03-w03-p03-056-finite-round} and the sentences of rank at most
\(n\) in the same way.

For the assertion about the Spoiler, suppose the Duplicator has no winning
strategy in the eternal game. By \cref{thm:p03-w03-p03-056-eternal} there is
an infinitary formula true on one side and false on the other at the initial
position. Maintain such a formula with its free labels pebbled. At a
negation exchange the two sides; at a conjunction choose a conjunct false
on the false side, and at a disjunction choose a disjunct true on the true
side. For \(\exists v_i\), play a witness on the true side; every answer
on the false side leaves the body true on one side and false on the other.
For \(\forall v_i\), play a counterexample on the false side and argue
dually. Reuse pebble \(i\), or first remove it in the separate-removal
variant. Every reduction passes to an immediate subformula, so its syntax
rank strictly decreases. There is no infinite descending sequence of
ordinals; hence each play reaches an atomic disagreement after finitely
many reductions, and therefore finitely many moves. By
\cref{thm:p03-w03-p03-056-atomic} this is a bad position. The finite length
may depend on the opponent's answers: the construction supplies no uniform
finite bound when the distinguishing formula has infinite rank.
\end{proof}
\end{annexitem}

\begin{annexitem}{annex:p03-w03-p03-056-definability}{thm:p03-w03-p03-056-definability}{Finite-variable definability on classes of finite structures}
\begin{proof}
\emph{(1)\(\Rightarrow\)(2).} Let \(\theta\) define \(Q\) on \(\mathcal C\).
By \cref{def:p03-w03-p03-056-grammar}, \(\theta\in\mathrm{L}^k_{\infty\omega}\)
for some \(k\geq1\). Let \(\mathbf A,\mathbf B\in\mathcal C\) with
\(\mathbf A\in Q\), and suppose the Duplicator wins
\(\mathrm{PG}^\infty_k(\mathbf A,\mathbf B)\). Then
\(\mathbf A\models\theta\), and by \cref{thm:p03-w03-p03-056-eternal}
(items~1 and~4, \(D=\emptyset\)) \(\mathbf B\models\theta\), so
\(\mathbf B\in Q\). Finiteness of the structures was not used.

\emph{(2)\(\Rightarrow\)(1).} Fix \(k\) as in item~2. Let \(S_k\) be the set of
\(\mathrm{FO}^k_{\omega\omega}(\sigma)\)-sentences; it is a set. Indeed, let \(W_0\) be the set of
atomic formulas and \(W_{j+1}\) the set consisting of \(W_j\) and the formulas obtained from members of
\(W_j\) by one application of (G2), of (G3) to a finite subset of \(W_j\), or
of (G4); each \(W_j\) is a set by recursion, and every
\(\mathrm{FO}^k_{\omega\omega}\)-formula lies in some \(W_j\), by induction on
its formation, since its formation tree is finite. So \(S_k\) is a subset of
\(\bigcup_jW_j\). For a
structure \(\mathbf A\) put
\(T(\mathbf A)=\{\varphi\in S_k:\mathbf A\models\varphi\}\subseteq S_k\), and
let
\[
 \mathcal T=\{T(\mathbf A):\mathbf A\in Q\},
\]
a subset of the power set of \(S_k\) (by separation), hence a set, though
\(Q\) itself may be a proper class. Put
\[
 \theta=\bigvee\Bigl\{\textstyle\bigwedge T:T\in\mathcal T\Bigr\}.
\]
Each \(\bigwedge T\) is formed by (G3) from a set of formulas with variables
among \(v_1,\ldots,v_k\) and no free labels, and so is \(\theta\); thus
\(\theta\) is an \(\mathrm{L}^k_{\infty\omega}\)-sentence.

Let \(\mathbf B\in\mathcal C\). If \(\mathbf B\in Q\), then
\(T(\mathbf B)\in\mathcal T\) and \(\mathbf B\models\bigwedge T(\mathbf B)\)
by definition of \(T(\mathbf B)\), so \(\mathbf B\models\theta\). Conversely
suppose \(\mathbf B\models\theta\). Then \(\mathbf B\models\bigwedge
T(\mathbf A)\) for some \(\mathbf A\in Q\), that is, every
\(\mathrm{FO}^k_{\omega\omega}\)-sentence true in \(\mathbf A\) is true in
\(\mathbf B\). If \(\varphi\in S_k\) is false in \(\mathbf A\), then
\(\neg\varphi\in S_k\) is true in \(\mathbf A\), hence in \(\mathbf B\), so
\(\varphi\) is false in \(\mathbf B\). Therefore
\(\mathbf A\equiv^k_{\omega\omega}\mathbf B\). Both structures are finite, so
the last clause of \cref{thm:p03-w03-p03-056-eternal} (with \(D=\emptyset\))
gives a winning strategy for the Duplicator in
\(\mathrm{PG}^\infty_k(\mathbf A,\mathbf B)\). As \(\mathbf A\in Q\) and
\(\mathbf A,\mathbf B\in\mathcal C\), item~2 gives \(\mathbf B\in Q\). Hence
\(\theta\) defines \(Q\) on \(\mathcal C\), and \(\theta\) has the form
stated after the theorem. (If \(Q\) is empty, \(\mathcal T=\emptyset\) and
\(\theta=\bigvee\emptyset\) is false in every structure, as required.)

\emph{Completeness of the method of \cref{def:p03-w03-p03-056-method}.} If
\(Q\) is not definable on \(\mathcal C\) by an
\(\mathrm{L}^\omega_{\infty\omega}\)-sentence, then item~2 fails for every
\(k\): for every \(k\geq1\) there are \(\mathbf A_k\in Q\) and
\(\mathbf B_k\in\mathcal C\setminus Q\) such that the Duplicator wins
\(\mathrm{PG}^\infty_k(\mathbf A_k,\mathbf B_k)\).
\end{proof}
The isomorphism-closure hypothesis on \(Q\) is not used in the proof; it
holds automatically when item~1 or item~2 holds, since isomorphic structures
satisfy the same sentences and the Duplicator wins every pebble game on
isomorphic structures (answer by the isomorphism; every position is then a
restriction of it, together with the constants).
\end{annexitem}

\begin{annexitem}{annex:p03-w03-p03-056-cliques}{thm:p03-w03-p03-056-cliques}{Cliques and pebble numbers}
\begin{proof}
The vocabulary \(\{E\}\) has no constants, so by
\cref{thm:p03-w03-p03-056-atomic} a position \((\alpha,\beta)\) with domain
\(D\) is good if and only if it agrees on the atoms \(v_i=v_j\) and
\(E(v_i,v_j)\), \(i,j\in D\). In a complete graph \(E(x,y)\) holds exactly
when \(x\neq y\). So the position is good if and only if
\[
 \alpha(i)=\alpha(j)\iff\beta(i)=\beta(j)\qquad(i,j\in D).\tag{\(\ast\)}
\]

(1) Let \(\mathbf A=\mathbf K_k\) and \(\mathbf B=\mathbf K_{k+1}\). The
empty position satisfies \((\ast)\). Suppose \((\ast)\) holds and the
Spoiler places pebble \(i\) on \(x\) in one structure; write \(\gamma\) for
the assignment of that structure, \(\gamma'\) for the other, and
\(D_0=D\setminus\{i\}\), so \(\lvert D_0\rvert\leq k-1\). If
\(x=\gamma(j_0)\) for some \(j_0\in D_0\), the Duplicator answers
\(y=\gamma'(j_0)\); then for \(j\in D_0\),
\(x=\gamma(j)\iff\gamma(j_0)=\gamma(j)\iff\gamma'(j_0)=\gamma'(j)\iff
y=\gamma'(j)\) by \((\ast)\). Otherwise she answers with an element \(y\) of
the other structure outside \(\gamma'[D_0]\); it exists because both
structures have at least \(k>\lvert D_0\rvert\) elements, and then \(x\) and
\(y\) both differ from every pebbled element with a label in \(D_0\). In
either case \((\ast)\) holds for the new position, the pairs with labels in
\(D_0\) being unchanged. A removal deletes a label and preserves \((\ast)\).
Hence every position is good, the Duplicator wins
\(\mathrm{PG}^\infty_k(\mathbf K_k,\mathbf K_{k+1})\), and
\(\mathbf K_k\equiv^k_{\infty\omega}\mathbf K_{k+1}\) by
\cref{thm:p03-w03-p03-056-eternal}.

(2) In rounds \(1,\ldots,k+1\) of
\(\mathrm{PG}^{k+1}_{k+1}(\mathbf K_k,\mathbf K_{k+1})\) let the Spoiler place
pebble \(j\) in round \(j\) on the \(j\)-th vertex of \(\mathbf K_{k+1}\),
with respect to a fixed enumeration. If some Duplicator strategy kept every
position good, after round \(k+1\) the position would have domain
\([k+1]\), with \(\beta\) injective, and by \((\ast)\) also \(\alpha\)
injective from \([k+1]\) into the \(k\)-element universe of \(\mathbf K_k\),
which is impossible. So no strategy is winning. The displayed sentence
\(\exists v_1\cdots\exists v_{k+1}\bigwedge\{v_i\neq v_j:i<j\}\) uses the
variables \(v_1,\ldots,v_{k+1}\) and finite conjunctions, so it is in
\(\mathrm{FO}^{k+1}_{\omega\omega}\) with rank \(k+1\). It holds in a
structure exactly when there are \(k+1\) pairwise distinct elements, so it
holds in \(\mathbf K_{k+1}\) and fails in \(\mathbf K_k\).
\end{proof}
\end{annexitem}

\begin{annexitem}{annex:p03-w03-p03-056-even}{thm:p03-w03-p03-056-even}{Even size is not finite-variable definable}
\begin{proof}
Suppose \(\theta\) were such a sentence. By
\cref{def:p03-w03-p03-056-grammar}, \(\theta\in\mathrm{L}^k_{\infty\omega}\)
for some \(k\geq1\). By \cref{thm:p03-w03-p03-056-cliques}(1),
\(\mathbf K_k\equiv^k_{\infty\omega}\mathbf K_{k+1}\), so \(\theta\) has the
same truth value in the two finite graphs. Exactly one of \(k\) and \(k+1\) is
even, so \(\theta\) is true in one graph of even size and one of odd size, or
false in both of them, contradicting the assumption in either case.
\end{proof}
\end{annexitem}

\begin{annexitem}{annex:p03-w03-p03-056-clique-query}{thm:p03-w03-p03-056-clique-query}{Clique queries need many variables}
\begin{proof}
\(\mathbf K_k\) and \(\mathbf K_{k+1}\) are finite graphs. In
\(\mathbf K_{k+1}\) all \(k+1\) vertices are pairwise adjacent; in
\(\mathbf K_k\) there are only \(k\) vertices, so no \(k+1\) pairwise
adjacent ones (adjacent vertices are distinct, \(E\) being irreflexive).
By \cref{thm:p03-w03-p03-056-cliques}(1),
\(\mathbf K_k\equiv^k_{\infty\omega}\mathbf K_{k+1}\), so any
\(\mathrm{L}^k_{\infty\omega}\)-sentence has the same truth value in both and
cannot define the query on finite graphs. The displayed sentence uses the
variables \(v_1,\ldots,v_{k+1}\) and a finite conjunction, so it is in
\(\mathrm{FO}^{k+1}_{\omega\omega}\). Under an assignment it holds exactly
when there are \(a_1,\ldots,a_{k+1}\) with \(E(a_i,a_j)\) for all \(i<j\);
since \(E\) is symmetric this gives \(E(a_i,a_j)\) for all \(i\neq j\), and
irreflexivity makes the \(a_i\) pairwise distinct. So it holds exactly when
there are \(k+1\) pairwise adjacent vertices.
\end{proof}
\end{annexitem}

\begin{annexitem}{annex:p03-w03-p03-056-orders}{thm:p03-w03-p03-056-orders}{Two pebbles count the length of a finite linear order}
\begin{proof}
Let \(\mathbf A=\mathbf L_m\) and \(\mathbf B=\mathbf L_n\) with \(m<n\). The
Spoiler plays only in \(\mathbf B\): in round \(j\) (\(1\leq j\leq n\)) he
places pebble \(\ell_j\) on the element \(j\), where \(\ell_j=1\) for odd
\(j\) and \(\ell_j=2\) for even \(j\). Let \(a_j\in A\) be the Duplicator's
answer in round \(j\). Since \(\ell_j\neq\ell_{j-1}\), after round
\(j\geq2\) the position contains the pairs \((a_{j-1},j-1)\) and
\((a_j,j)\), pebble \(\ell_{j-1}\) not having moved in round \(j\). Suppose
some Duplicator strategy kept every position good for \(n\) rounds. Then
after round \(j\geq2\) the map \(\rho\) is a partial isomorphism
containing these two pairs, and since \(j-1<j\) in \(\mathbf B\), it gives
\(a_{j-1}<a_j\). Thus \(a_1<a_2<\cdots<a_n\) are \(n\) distinct elements of
the \(m\)-element set \(\{1,\ldots,m\}\), which is impossible as \(m<n\). So
the Duplicator has no winning strategy in
\(\mathrm{PG}^n_2(\mathbf L_m,\mathbf L_n)\). By
\cref{thm:p03-w03-p03-056-finite-round} with \(D=\emptyset\), items~1 and~4,
\(\mathbf L_m\not\equiv^{2,n}_{\omega\omega}\mathbf L_n\): some
\(\mathrm{FO}^2_{\omega\omega}\)-sentence of rank at most \(n\) holds in
exactly one of them.
\end{proof}
\end{annexitem}

\begin{annexitem}{annex:p03-w03-p03-056-successor}{thm:p03-w03-p03-056-successor}{First-order two-variable equivalence without infinitary equivalence}
\begin{proof}
\emph{Coordinates.} Regard \(\mathbf A\) as having one component, with
coordinate \(x\) for the element \(x\); in \(\mathbf B\) the element
\((e,x)\) lies in component \(e\) and has coordinate \(x\). In both
structures \(E(u,u')\) holds exactly when \(u,u'\) are in the same component
and the coordinate of \(u'\) is one more than that of \(u\). For \(u\) in a
component and an integer \(z\), write \(u+z\) for the element of the same
component whose coordinate is that of \(u\) plus \(z\), when it exists. The
\emph{depth} \(d(u)\) is the coordinate of \(u\) if \(u\) is in \(\mathbf A\)
or in component \(0\) of \(\mathbf B\), and \(\infty\) if \(u\) is in
component \(1\); it is the largest \(z\) for which \(u-z\) exists. For
\(u,u'\) in one structure, \(\Delta(u,u')\) is the coordinate of \(u'\) minus
that of \(u\) when they are in the same component, and is undefined
otherwise. For \(t\geq1\), the pair \((u,u')\) is \(t\)-\emph{near} when
\(\Delta(u,u')\) is defined and \(\lvert\Delta(u,u')\rvert<t\).

\emph{Invariant.} Put \(t_r=2^{r+1}\), so \(t_{r+1}=2t_r\) and \(t_0=2\). A
position \((\alpha,\beta)\) with domain \(D\subseteq\{1,2\}\) satisfies
\(\mathrm{Inv}_r\) when
\begin{enumerate}
\item[(I1)] for each \(i\in D\), \(d(\alpha(i))=d(\beta(i))\) or both are
      \(\geq t_r\); and
\item[(I2)] if \(D=\{1,2\}\), then \((\alpha(1),\alpha(2))\) is
      \(t_r\)-near if and only if \((\beta(1),\beta(2))\) is, and in that
      case \(\Delta(\alpha(1),\alpha(2))=\Delta(\beta(1),\beta(2))\).
\end{enumerate}

\emph{Claim A: \(\mathrm{Inv}_r\) implies good.} There are no constants, so
by \cref{thm:p03-w03-p03-056-atomic} it suffices to check the atoms
\(v_i=v_j\) and \(E(v_i,v_j)\) for \(i,j\in D\). For \(i=j\): \(v_i=v_i\) is
true and \(E(v_i,v_i)\) is false in both structures, since no element is
its own \(E\)-successor. For \(\{i,j\}=\{1,2\}\): \(\alpha(i)=\alpha(j)\)
holds iff \(\Delta(\alpha(i),\alpha(j))=0\), and \(E(\alpha(i),\alpha(j))\)
holds iff \(\Delta(\alpha(i),\alpha(j))=1\). In either case the pair is
\(t_r\)-near, because \(t_r\geq2\), so by (I2) the \(\beta\)-pair is near with
the same \(\Delta\) and the atom holds in \(\mathbf B\). The converse is the
same argument from \(\mathbf B\). (Note \(\Delta(u',u)=-\Delta(u,u')\), so the
order of the labels does not matter.)

\emph{Claim B: \(\mathrm{Inv}_{r+1}\) implies \(\mathrm{Inv}_r\).} For (I1),
depths \(\geq t_{r+1}\) are \(\geq t_r\). For (I2), a \(t_r\)-near pair is
\(t_{r+1}\)-near, so its partner is \(t_{r+1}\)-near with the same
\(\Delta\), whose absolute value is \(<t_r\); so the partner is
\(t_r\)-near. Both directions are symmetric. Since \(\mathrm{Inv}_r\) for a
position implies \(\mathrm{Inv}_r\) for its restrictions, a removal from an
\(\mathrm{Inv}_{r+1}\)-position gives an \(\mathrm{Inv}_r\)-position.

\emph{Claim C: answering placements.} Let \((\alpha,\beta)\) satisfy
\(\mathrm{Inv}_{r+1}\) and let the Spoiler place pebble \(i\). Let \(j\) be
the other label. The pair with label \(i\), if present, is lifted, and
\(\mathrm{Inv}_r\) for the pair \(j\) (if \(j\in D\)) follows from
\(\mathrm{Inv}_{r+1}\) by Claim B. It remains to answer so that the new pair
\((x,y)\), \(x\in A\), \(y\in B\), satisfies (I1) for \(t_r\), and, if
\(j\in D\), satisfies (I2) for \(t_r\) together with the pair
\((a,b)=(\alpha(j),\beta(j))\). Recall from \(\mathrm{Inv}_{r+1}\) that
\(d(a)=d(b)\) or both are \(\geq t_{r+1}\).

\emph{Case \(j\notin D\).} If the Spoiler chose \(x\in A\): answer
\(y=(0,x)\) when \(x<t_r\), so the depths are equal, and \(y=(1,0)\)
otherwise, of depth \(\infty\geq t_r\). If he chose \(y\in B\): answer with
the coordinate of \(y\) when \(d(y)<t_r\) (then \(y\) is in component \(0\)
and the depths are equal), and with \(x=t_r\) otherwise.

\emph{Case \(j\in D\), Spoiler chooses \(x\in A\).}
\begin{itemize}
\item \(\lvert x-a\rvert<t_r\). Answer \(y=b+(x-a)\). It exists: if
      \(x\geq a\) because components are unbounded above; if \(x<a\)
      because \(a-x\leq a=d(a)\) and either \(d(b)=d(a)\) or
      \(d(b)\geq t_{r+1}>t_r>a-x\). Then \(\Delta(b,y)=x-a=\Delta(a,x)\),
      both pairs are \(t_r\)-near, and (I2) holds. For (I1):
      \(d(x)=d(a)+(x-a)\) and \(d(y)=d(b)+(x-a)\) (with
      \(\infty+z=\infty\)); if \(d(a)=d(b)\) these are equal, and otherwise
      \(d(a),d(b)\geq t_{r+1}\) gives \(d(x),d(y)\geq t_{r+1}-t_r=t_r\).
\item \(\lvert x-a\rvert\geq t_r\) and \(x<t_r\). Answer \(y=(0,x)\); the
      depths are equal. The pair \((b,y)\) is not \(t_r\)-near: if \(b\) is
      in component \(1\), \(\Delta\) is undefined; if \(b=(0,b')\) and
      \(d(a)=d(b)\), then \(b'=a\) and \(\lvert\Delta(b,y)\rvert=
      \lvert x-a\rvert\geq t_r\); if \(b=(0,b')\) and
      \(d(a),d(b)\geq t_{r+1}\), then \(b'-x>t_{r+1}-t_r=t_r\).
\item \(\lvert x-a\rvert\geq t_r\) and \(x\geq t_r\). Answer \(y=(1,0)\) if
      \(b\) is in component \(0\), and \(y=(0,t_r)\) if \(b\) is in component
      \(1\). In both cases \(d(y)\geq t_r\) and \(y\) lies in the other
      component from \(b\), so \((b,y)\) is not near, and neither is
      \((a,x)\).
\end{itemize}

\emph{Case \(j\in D\), Spoiler chooses \(y\in B\).}
\begin{itemize}
\item \((b,y)\) is \(t_r\)-near. Answer \(x=a+\Delta(b,y)\). It exists: if
      \(\Delta(b,y)\geq0\) trivially; if \(\Delta(b,y)<0\) and
      \(d(a)=d(b)\), then \(b=(0,a)\), \(y\) is in component \(0\) with
      coordinate \(a+\Delta(b,y)\geq0\); if \(d(a)\geq t_{r+1}\), then
      \(a\geq t_{r+1}>t_r>-\Delta(b,y)\). Both pairs are near with the same
      \(\Delta\), and (I1) follows exactly as in the first bullet above.
\item \((b,y)\) is not \(t_r\)-near and \(d(y)<t_r\). Then \(y=(0,y')\) with
      \(y'<t_r\); answer \(x=y'\), so the depths are equal. If
      \(d(a)=d(b)\), then \(b=(0,a)\), \(\Delta(b,y)=y'-a\) is defined, and
      not being near means \(\lvert x-a\rvert\geq t_r\). Otherwise
      \(a=d(a)\geq t_{r+1}\) and \(a-x>t_{r+1}-t_r=t_r\). Either way
      \((a,x)\) is not near.
\item \((b,y)\) is not \(t_r\)-near and \(d(y)\geq t_r\). Answer
      \(x=a+t_r\): then \(d(x)=a+t_r\geq t_r\) and
      \(\lvert\Delta(a,x)\rvert=t_r\), so \((a,x)\) is not near.
\end{itemize}
In every case (I1) and (I2) hold for \(t_r\), so the new position satisfies
\(\mathrm{Inv}_r\).

\emph{Item 1.} Fix \(n\). The empty position satisfies \(\mathrm{Inv}_n\).
By Claims B and C, if the position before a round satisfies
\(\mathrm{Inv}_{r+1}\), the Duplicator can make the position after it
satisfy \(\mathrm{Inv}_r\). Hence after \(j\) rounds the position satisfies
\(\mathrm{Inv}_{n-j}\), for \(0\leq j\leq n\), and every position is good by
Claim A. So the Duplicator wins \(\mathrm{PG}^n_2(\mathbf A,\mathbf B)\). Every
\(\mathrm{FO}^2_{\omega\omega}\)-sentence has some finite rank \(n\), and
\cref{thm:p03-w03-p03-056-finite-round} (items~1 and~4) gives agreement on
it; so \(\mathbf A\equiv^2_{\omega\omega}\mathbf B\).

\emph{Item 2.} Each \(\delta_j\) is built by (G1)--(G4) from \(v_1,v_2\), and
\(\operatorname{fr}(\delta_j)=\{1\}\): for \(\delta_0\) directly, and if
\(\operatorname{fr}(\delta_j)=\{1\}\) then
\(\exists v_1(v_1=v_2\land\delta_j)\) has free labels \(\{2\}\), the
conjunction with \(E(v_2,v_1)\) has \(\{1,2\}\), and \(\exists v_2\) leaves
\(\{1\}\). By induction on \(j\), \(\delta_j\) holds of an element \(u\) (in
either structure) iff \(u-j\) exists, that is, iff \(d(u)\geq j\): for
\(j=0\) both are always true; and \(\delta_{j+1}\) holds of \(u\) iff some
\(w\) with \(E(w,u)\) satisfies \(\delta_j\), the inner quantifier
rebinding \(v_1\) to the value of \(v_2\) as in
\cref{annex:p03-w03-p03-056-walks}. The only candidate is \(w=u-1\), so this
holds iff \(u-1\) exists and \(d(u-1)\geq j\), iff \(d(u)\geq j+1\).
Hence \(\exists v_1\bigwedge\{\delta_j:j\in\N\}\) holds iff some element has
depth \(\infty\). It holds in \(\mathbf B\), at \((1,0)\), and fails in
\(\mathbf A\), where every depth is finite. So
\(\mathbf A\not\equiv^2_{\infty\omega}\mathbf B\), and by
\cref{thm:p03-w03-p03-056-eternal} the Duplicator has no winning strategy in
\(\mathrm{PG}^\infty_2(\mathbf A,\mathbf B)\).
\end{proof}
\end{annexitem}

\begin{annexitem}{annex:p03-w03-p03-056-ef-bridge}{thm:p03-w03-p03-056-ef-bridge}{Pebble and Ehrenfeucht--Fra\"iss\'e games}
\begin{proof}
For a constant-free vocabulary a position is good exactly when its set of
pairs is a finite partial isomorphism in the sense of
\cref{def:p03-ef-game}, which also requires respect for equality and the
preservation and reflection of relations.

\emph{From the Ehrenfeucht--Fra\"iss\'e game to the pebble game.} Let
\(\Sigma\) be a winning strategy in the \(n\)-round
Ehrenfeucht--Fra\"iss\'e game. In \(\mathrm{PG}^n_n\) the Duplicator keeps a
record of an Ehrenfeucht--Fra\"iss\'e play: each placement move (of any
label on any element) is entered as the next Ehrenfeucht--Fra\"iss\'e move
of the Spoiler on the same element, and she answers with \(\Sigma\)'s
answer; removals enter nothing. At most \(n\) moves are recorded. Any
recorded play can be continued to a full \(n\)-round play following
\(\Sigma\), whose pairs form a partial isomorphism; hence the recorded pairs
form a partial isomorphism \(h\) (a restriction of that one). Each pair now
on the board was created by a recorded placement, so the board's pairs form
a subset of the graph of \(h\), and a restriction of a partial isomorphism
to a subset of its domain is a partial isomorphism. Every position is good.

\emph{From the pebble game to the Ehrenfeucht--Fra\"iss\'e game.} Let
\(\Sigma\) be winning in \(\mathrm{PG}^n_n\). In round \(j\) of the
Ehrenfeucht--Fra\"iss\'e game the Duplicator treats the Spoiler's choice as
the placement of pebble \(j\), which is not yet on the board, on the same
element, and answers as \(\Sigma\) does. After \(n\) rounds the pebble
position has domain \([n]\) and its pairs are exactly the recorded
Ehrenfeucht--Fra\"iss\'e pairs \(a_j\mapsto b_j\); this position is good, so
the pairing is a finite partial isomorphism and the Duplicator wins.
\end{proof}
\end{annexitem}
